% ---------------------------------------------------------------------------
% Sections 5 (Discussion and Limitations, five short paragraphs) and 6
% (Conclusion, three sentences) of the 2026-10-07 draft.  % ---------------------------------------------------------------------------
% Sections 5 (Discussion and Limitations, five short paragraphs) and 6
% (Conclusion, three sentences) of the 2026-10-07 draft.  % ---------------------------------------------------------------------------
% Sections 5 (Discussion and Limitations, five short paragraphs) and 6
% (Conclusion, three sentences) of the 2026-10-07 draft.  % ---------------------------------------------------------------------------
% Sections 5 (Discussion and Limitations, five short paragraphs) and 6
% (Conclusion, three sentences) of the 2026-10-07 draft.  \input{discussion}
% after results.tex.  Claims here are the ones Sec. 4 supports on disk; the
% ones that wait on the FieldProc reruns or E10 are marked \todo.
% ---------------------------------------------------------------------------
\providecommand{\CropCode}{\textsc{CropCode}}
\providecommand{\FieldProc}{\textsc{FieldProc}}
\providecommand{\todo}[1]{\textcolor{red}{#1}}

\section{Discussion and Limitations}
\label{sec:discussion}

\paragraph{Why a stream that holds the answer hurts.}
The parameter stream is the only input that states the phenotype verbatim, and
it is the one input never available in the field. Every arm in which the encoder
saw the parameter stream learned a worse sensor representation, and on the size
traits the loss grew with the stream's visibility (Tab.~\ref{tab:control}).
Even as a pure decoder target, the parameter stream lowered roll, leaf length
and leaf width, the traits stored per leaf in the parameter tokens. One
explanation, consistent with modality competition and negative transfer, is that
a privileged stream lets the encoder route a trait through tokens that are
absent at test time, and the reconstruction pressure of a per-leaf target then
shifts the latent toward the plant's size, which is already the direction of its
first principal component. The ablations support a practical rule: a generative
modality belongs in the data and in the targets, not in the encoder input. Our
headline model does not yet follow this rule: it was pretrained with the
parameters as an encoder input, and its distillation used the same model as
teacher, so its frozen representation corresponds to the four-stream arm. A
headline run at $\rho_\theta{=}1$ would be the direct test. We expect the rule
to transfer to any multimodal model with access to a program, a CAD sketch, or a
simulator state, and the gated ratio $\rho_\theta$ is a one-line way to test
it.

\paragraph{Which metric decides.}
Chamfer and the frozen probe agree on the coarse ranking of arms and disagree on
the design decisions (Fig.~\ref{fig:metrics}): depth, more training plants, and
the parameter stream each move Chamfer by tens of percent and leave the probe
nearly unchanged, while a partial cloud from one camera raises Chamfer by 2.6 to
6.5 times but changes the probe by at most $0.03$. On slender geometry, Chamfer
measures surface proximity and the probe measures recovered structure, so we use
the probe to make decisions and report Chamfer as the decoder's surface error.
The probe has its own limitations: it requires labels and a held-out split, and
its $R^2$ is comparable between models on one split but is not an absolute
measure across datasets (our held-out plants are extreme-enriched). Sorghum's
four named traits also reduce to two generator variables, so our claims about
independent traits rely on maize. On the training side, we added the Sinkhorn
term because Chamfer barely penalizes a missing blade; its ablation is pending
(\todo{E7}).

\paragraph{From complete clouds to a camera.}
Every sorghum model was pretrained on complete clouds of isolated plants, and
Sec.~\ref{sec:regime2} measures the effect of partial clouds. \todo{Maize is the other
case: a maize \texttt{pointcloud\_cam.ply} is the surface visible to its own camera
(2\,\% of a cloud's points hidden from that camera, 25\,\% from another; the
depth image back-projects to within 1\,mm), so every maize model trained on
single-camera clouds. The maize study inverts the question (what three and ten
fused cameras add; \texttt{eval/pc\_view\_eval\_maize.py}) and is queued.} The phenotype information survives a
single camera's view, but the latent shifts with coverage, so a readout
calibrated on complete scans is wrong on partial ones and must be refit. The
decoder reproduces what it is given and does not hallucinate the hidden half.
Viewpoint, not morphology, determines the difficulty of inverting a photograph.
The recipe tolerates partial views best: the recovered insertion heights of one
plant vary by a tenth of their spread across its ten photographs, and a
re-rendered recipe is complete even where the cloud is not. In both species,
training on rows removes most of the neighbor penalty (91\,\% in maize and
85\,\% in sorghum, Sec.~\ref{sec:regime3}). The real scans
(Sec.~\ref{sec:res-field}) are not yet evaluated; until they are, the
one-camera transfer gap is our prediction for the sim-to-real penalty. A generator also contains unintended
regularities, such as the fixed leaf azimuth that gave a pose-aware baseline a
canonical orientation, and benchmark users should expect more of them.

\paragraph{What distillation does and does not buy.}
Distillation from a full-modality teacher improves every single-sensor student
and concentrates the gain on the weakest sources and the hardest plants
(Tab.~\ref{tab:anysubset}, Fig.~\ref{fig:views}B), and supervising all four
targets beats a parameter specialist at parameter regression. Distillation does
not make the model view-robust: the elevation curve moves down and keeps its
shape, so view robustness requires a different method, most plausibly treating
views as an explicit augmentation axis during pretraining. Distillation also
showed that an unpredictable field is not neutral: the three waviness parameters
of the first sorghum release stayed at zero skill and became more view-dependent
under distillation, and the final generator does not include them. A
parameterization meant to be regressed should be restricted to the fields that
the sensors can determine.

\paragraph{Scope and compute.}
The ablations use one seed per arm. The run-to-run floor comes from a single
near-duplicate pair (2\,\% in Chamfer, at most $0.011$ in $R^2$), and the maize
scaling differences, although significant under the paired bootstrap, could
still be run-to-run variance (Sec.~\ref{sec:res-repr}). Two generators and two
species-specific models are studied, and the 1{,}000-epoch headline model is
trained once per species. The procedural prior that makes inversion tractable
also limits it: a plant that the generator cannot express cannot be recovered,
and an organ class that the generator renders without variation (tassel droop in
maize is below zero for every encoder of this design) limits the data as much as
the model. Real scans and the \FieldProc{} release at
10K--100K plants remain \todo{in progress}.

\section{Conclusion}
\label{sec:conclusion}

\CropCode{} recovers a complete, editable procedural plant from the data of
whichever sensor a field campaign used, in one forward pass, by using the
generator for both the training data and the output representation. Two
findings from its ablations apply beyond plants: the generator's parameters are
valuable targets but harmful encoder inputs, and on slender geometry a
downstream probe, not Chamfer distance, should decide what a modality is worth.
A single camera's partial view requires refitting the readout and leaves a
measured transfer gap; the row-occluded and real-scan evaluations will show
whether the generator's prior covers the remaining difference to field data.
\todo{Sharpen once E10 lands.}

% after results.tex.  Claims here are the ones Sec. 4 supports on disk; the
% ones that wait on the FieldProc reruns or E10 are marked \todo.
% ---------------------------------------------------------------------------
\providecommand{\CropCode}{\textsc{CropCode}}
\providecommand{\FieldProc}{\textsc{FieldProc}}
\providecommand{\todo}[1]{\textcolor{red}{#1}}

\section{Discussion and Limitations}
\label{sec:discussion}

\paragraph{Why a stream that holds the answer hurts.}
The parameter stream is the only input that states the phenotype verbatim, and
it is the one input never available in the field. Every arm in which the encoder
saw the parameter stream learned a worse sensor representation, and on the size
traits the loss grew with the stream's visibility (Tab.~\ref{tab:control}).
Even as a pure decoder target, the parameter stream lowered roll, leaf length
and leaf width, the traits stored per leaf in the parameter tokens. One
explanation, consistent with modality competition and negative transfer, is that
a privileged stream lets the encoder route a trait through tokens that are
absent at test time, and the reconstruction pressure of a per-leaf target then
shifts the latent toward the plant's size, which is already the direction of its
first principal component. The ablations support a practical rule: a generative
modality belongs in the data and in the targets, not in the encoder input. Our
headline model does not yet follow this rule: it was pretrained with the
parameters as an encoder input, and its distillation used the same model as
teacher, so its frozen representation corresponds to the four-stream arm. A
headline run at $\rho_\theta{=}1$ would be the direct test. We expect the rule
to transfer to any multimodal model with access to a program, a CAD sketch, or a
simulator state, and the gated ratio $\rho_\theta$ is a one-line way to test
it.

\paragraph{Which metric decides.}
Chamfer and the frozen probe agree on the coarse ranking of arms and disagree on
the design decisions (Fig.~\ref{fig:metrics}): depth, more training plants, and
the parameter stream each move Chamfer by tens of percent and leave the probe
nearly unchanged, while a partial cloud from one camera raises Chamfer by 2.6 to
6.5 times but changes the probe by at most $0.03$. On slender geometry, Chamfer
measures surface proximity and the probe measures recovered structure, so we use
the probe to make decisions and report Chamfer as the decoder's surface error.
The probe has its own limitations: it requires labels and a held-out split, and
its $R^2$ is comparable between models on one split but is not an absolute
measure across datasets (our held-out plants are extreme-enriched). Sorghum's
four named traits also reduce to two generator variables, so our claims about
independent traits rely on maize. On the training side, we added the Sinkhorn
term because Chamfer barely penalizes a missing blade; its ablation is pending
(\todo{E7}).

\paragraph{From complete clouds to a camera.}
Every sorghum model was pretrained on complete clouds of isolated plants, and
Sec.~\ref{sec:regime2} measures the effect of partial clouds. \todo{Maize is the other
case: a maize \texttt{pointcloud\_cam.ply} is the surface visible to its own camera
(2\,\% of a cloud's points hidden from that camera, 25\,\% from another; the
depth image back-projects to within 1\,mm), so every maize model trained on
single-camera clouds. The maize study inverts the question (what three and ten
fused cameras add; \texttt{eval/pc\_view\_eval\_maize.py}) and is queued.} The phenotype information survives a
single camera's view, but the latent shifts with coverage, so a readout
calibrated on complete scans is wrong on partial ones and must be refit. The
decoder reproduces what it is given and does not hallucinate the hidden half.
Viewpoint, not morphology, determines the difficulty of inverting a photograph.
The recipe tolerates partial views best: the recovered insertion heights of one
plant vary by a tenth of their spread across its ten photographs, and a
re-rendered recipe is complete even where the cloud is not. In both species,
training on rows removes most of the neighbor penalty (91\,\% in maize and
85\,\% in sorghum, Sec.~\ref{sec:regime3}). The real scans
(Sec.~\ref{sec:res-field}) are not yet evaluated; until they are, the
one-camera transfer gap is our prediction for the sim-to-real penalty. A generator also contains unintended
regularities, such as the fixed leaf azimuth that gave a pose-aware baseline a
canonical orientation, and benchmark users should expect more of them.

\paragraph{What distillation does and does not buy.}
Distillation from a full-modality teacher improves every single-sensor student
and concentrates the gain on the weakest sources and the hardest plants
(Tab.~\ref{tab:anysubset}, Fig.~\ref{fig:views}B), and supervising all four
targets beats a parameter specialist at parameter regression. Distillation does
not make the model view-robust: the elevation curve moves down and keeps its
shape, so view robustness requires a different method, most plausibly treating
views as an explicit augmentation axis during pretraining. Distillation also
showed that an unpredictable field is not neutral: the three waviness parameters
of the first sorghum release stayed at zero skill and became more view-dependent
under distillation, and the final generator does not include them. A
parameterization meant to be regressed should be restricted to the fields that
the sensors can determine.

\paragraph{Scope and compute.}
The ablations use one seed per arm. The run-to-run floor comes from a single
near-duplicate pair (2\,\% in Chamfer, at most $0.011$ in $R^2$), and the maize
scaling differences, although significant under the paired bootstrap, could
still be run-to-run variance (Sec.~\ref{sec:res-repr}). Two generators and two
species-specific models are studied, and the 1{,}000-epoch headline model is
trained once per species. The procedural prior that makes inversion tractable
also limits it: a plant that the generator cannot express cannot be recovered,
and an organ class that the generator renders without variation (tassel droop in
maize is below zero for every encoder of this design) limits the data as much as
the model. Real scans and the \FieldProc{} release at
10K--100K plants remain \todo{in progress}.

\section{Conclusion}
\label{sec:conclusion}

\CropCode{} recovers a complete, editable procedural plant from the data of
whichever sensor a field campaign used, in one forward pass, by using the
generator for both the training data and the output representation. Two
findings from its ablations apply beyond plants: the generator's parameters are
valuable targets but harmful encoder inputs, and on slender geometry a
downstream probe, not Chamfer distance, should decide what a modality is worth.
A single camera's partial view requires refitting the readout and leaves a
measured transfer gap; the row-occluded and real-scan evaluations will show
whether the generator's prior covers the remaining difference to field data.
\todo{Sharpen once E10 lands.}

% after results.tex.  Claims here are the ones Sec. 4 supports on disk; the
% ones that wait on the FieldProc reruns or E10 are marked \todo.
% ---------------------------------------------------------------------------
\providecommand{\CropCode}{\textsc{CropCode}}
\providecommand{\FieldProc}{\textsc{FieldProc}}
\providecommand{\todo}[1]{\textcolor{red}{#1}}

\section{Discussion and Limitations}
\label{sec:discussion}

\paragraph{Why a stream that holds the answer hurts.}
The parameter stream is the only input that states the phenotype verbatim, and
it is the one input never available in the field. Every arm in which the encoder
saw the parameter stream learned a worse sensor representation, and on the size
traits the loss grew with the stream's visibility (Tab.~\ref{tab:control}).
Even as a pure decoder target, the parameter stream lowered roll, leaf length
and leaf width, the traits stored per leaf in the parameter tokens. One
explanation, consistent with modality competition and negative transfer, is that
a privileged stream lets the encoder route a trait through tokens that are
absent at test time, and the reconstruction pressure of a per-leaf target then
shifts the latent toward the plant's size, which is already the direction of its
first principal component. The ablations support a practical rule: a generative
modality belongs in the data and in the targets, not in the encoder input. Our
headline model does not yet follow this rule: it was pretrained with the
parameters as an encoder input, and its distillation used the same model as
teacher, so its frozen representation corresponds to the four-stream arm. A
headline run at $\rho_\theta{=}1$ would be the direct test. We expect the rule
to transfer to any multimodal model with access to a program, a CAD sketch, or a
simulator state, and the gated ratio $\rho_\theta$ is a one-line way to test
it.

\paragraph{Which metric decides.}
Chamfer and the frozen probe agree on the coarse ranking of arms and disagree on
the design decisions (Fig.~\ref{fig:metrics}): depth, more training plants, and
the parameter stream each move Chamfer by tens of percent and leave the probe
nearly unchanged, while a partial cloud from one camera raises Chamfer by 2.6 to
6.5 times but changes the probe by at most $0.03$. On slender geometry, Chamfer
measures surface proximity and the probe measures recovered structure, so we use
the probe to make decisions and report Chamfer as the decoder's surface error.
The probe has its own limitations: it requires labels and a held-out split, and
its $R^2$ is comparable between models on one split but is not an absolute
measure across datasets (our held-out plants are extreme-enriched). Sorghum's
four named traits also reduce to two generator variables, so our claims about
independent traits rely on maize. On the training side, we added the Sinkhorn
term because Chamfer barely penalizes a missing blade; its ablation is pending
(\todo{E7}).

\paragraph{From complete clouds to a camera.}
Every sorghum model was pretrained on complete clouds of isolated plants, and
Sec.~\ref{sec:regime2} measures the effect of partial clouds. \todo{Maize is the other
case: a maize \texttt{pointcloud\_cam.ply} is the surface visible to its own camera
(2\,\% of a cloud's points hidden from that camera, 25\,\% from another; the
depth image back-projects to within 1\,mm), so every maize model trained on
single-camera clouds. The maize study inverts the question (what three and ten
fused cameras add; \texttt{eval/pc\_view\_eval\_maize.py}) and is queued.} The phenotype information survives a
single camera's view, but the latent shifts with coverage, so a readout
calibrated on complete scans is wrong on partial ones and must be refit. The
decoder reproduces what it is given and does not hallucinate the hidden half.
Viewpoint, not morphology, determines the difficulty of inverting a photograph.
The recipe tolerates partial views best: the recovered insertion heights of one
plant vary by a tenth of their spread across its ten photographs, and a
re-rendered recipe is complete even where the cloud is not. In both species,
training on rows removes most of the neighbor penalty (91\,\% in maize and
85\,\% in sorghum, Sec.~\ref{sec:regime3}). The real scans
(Sec.~\ref{sec:res-field}) are not yet evaluated; until they are, the
one-camera transfer gap is our prediction for the sim-to-real penalty. A generator also contains unintended
regularities, such as the fixed leaf azimuth that gave a pose-aware baseline a
canonical orientation, and benchmark users should expect more of them.

\paragraph{What distillation does and does not buy.}
Distillation from a full-modality teacher improves every single-sensor student
and concentrates the gain on the weakest sources and the hardest plants
(Tab.~\ref{tab:anysubset}, Fig.~\ref{fig:views}B), and supervising all four
targets beats a parameter specialist at parameter regression. Distillation does
not make the model view-robust: the elevation curve moves down and keeps its
shape, so view robustness requires a different method, most plausibly treating
views as an explicit augmentation axis during pretraining. Distillation also
showed that an unpredictable field is not neutral: the three waviness parameters
of the first sorghum release stayed at zero skill and became more view-dependent
under distillation, and the final generator does not include them. A
parameterization meant to be regressed should be restricted to the fields that
the sensors can determine.

\paragraph{Scope and compute.}
The ablations use one seed per arm. The run-to-run floor comes from a single
near-duplicate pair (2\,\% in Chamfer, at most $0.011$ in $R^2$), and the maize
scaling differences, although significant under the paired bootstrap, could
still be run-to-run variance (Sec.~\ref{sec:res-repr}). Two generators and two
species-specific models are studied, and the 1{,}000-epoch headline model is
trained once per species. The procedural prior that makes inversion tractable
also limits it: a plant that the generator cannot express cannot be recovered,
and an organ class that the generator renders without variation (tassel droop in
maize is below zero for every encoder of this design) limits the data as much as
the model. Real scans and the \FieldProc{} release at
10K--100K plants remain \todo{in progress}.

\section{Conclusion}
\label{sec:conclusion}

\CropCode{} recovers a complete, editable procedural plant from the data of
whichever sensor a field campaign used, in one forward pass, by using the
generator for both the training data and the output representation. Two
findings from its ablations apply beyond plants: the generator's parameters are
valuable targets but harmful encoder inputs, and on slender geometry a
downstream probe, not Chamfer distance, should decide what a modality is worth.
A single camera's partial view requires refitting the readout and leaves a
measured transfer gap; the row-occluded and real-scan evaluations will show
whether the generator's prior covers the remaining difference to field data.
\todo{Sharpen once E10 lands.}

% after results.tex.  Claims here are the ones Sec. 4 supports on disk; the
% ones that wait on the FieldProc reruns or E10 are marked \todo.
% ---------------------------------------------------------------------------
\providecommand{\CropCode}{\textsc{CropCode}}
\providecommand{\FieldProc}{\textsc{FieldProc}}
\providecommand{\todo}[1]{\textcolor{red}{#1}}

\section{Discussion and Limitations}
\label{sec:discussion}

\paragraph{Why a stream that holds the answer hurts.}
The parameter stream is the only input that states the phenotype verbatim, and
it is the one input never available in the field. Every arm in which the encoder
saw the parameter stream learned a worse sensor representation, and on the size
traits the loss grew with the stream's visibility (Tab.~\ref{tab:control}).
Even as a pure decoder target, the parameter stream lowered roll, leaf length
and leaf width, the traits stored per leaf in the parameter tokens. One
explanation, consistent with modality competition and negative transfer, is that
a privileged stream lets the encoder route a trait through tokens that are
absent at test time, and the reconstruction pressure of a per-leaf target then
shifts the latent toward the plant's size, which is already the direction of its
first principal component. The ablations support a practical rule: a generative
modality belongs in the data and in the targets, not in the encoder input. Our
headline model does not yet follow this rule: it was pretrained with the
parameters as an encoder input, and its distillation used the same model as
teacher, so its frozen representation corresponds to the four-stream arm. A
headline run at $\rho_\theta{=}1$ would be the direct test. We expect the rule
to transfer to any multimodal model with access to a program, a CAD sketch, or a
simulator state, and the gated ratio $\rho_\theta$ is a one-line way to test
it.

\paragraph{Which metric decides.}
Chamfer and the frozen probe agree on the coarse ranking of arms and disagree on
the design decisions (Fig.~\ref{fig:metrics}): depth, more training plants, and
the parameter stream each move Chamfer by tens of percent and leave the probe
nearly unchanged, while a partial cloud from one camera raises Chamfer by 2.6 to
6.5 times but changes the probe by at most $0.03$. On slender geometry, Chamfer
measures surface proximity and the probe measures recovered structure, so we use
the probe to make decisions and report Chamfer as the decoder's surface error.
The probe has its own limitations: it requires labels and a held-out split, and
its $R^2$ is comparable between models on one split but is not an absolute
measure across datasets (our held-out plants are extreme-enriched). Sorghum's
four named traits also reduce to two generator variables, so our claims about
independent traits rely on maize. On the training side, we added the Sinkhorn
term because Chamfer barely penalizes a missing blade; its ablation is pending
(\todo{E7}).

\paragraph{From complete clouds to a camera.}
Every sorghum model was pretrained on complete clouds of isolated plants, and
Sec.~\ref{sec:regime2} measures the effect of partial clouds. \todo{Maize is the other
case: a maize \texttt{pointcloud\_cam.ply} is the surface visible to its own camera
(2\,\% of a cloud's points hidden from that camera, 25\,\% from another; the
depth image back-projects to within 1\,mm), so every maize model trained on
single-camera clouds. The maize study inverts the question (what three and ten
fused cameras add; \texttt{eval/pc\_view\_eval\_maize.py}) and is queued.} The phenotype information survives a
single camera's view, but the latent shifts with coverage, so a readout
calibrated on complete scans is wrong on partial ones and must be refit. The
decoder reproduces what it is given and does not hallucinate the hidden half.
Viewpoint, not morphology, determines the difficulty of inverting a photograph.
The recipe tolerates partial views best: the recovered insertion heights of one
plant vary by a tenth of their spread across its ten photographs, and a
re-rendered recipe is complete even where the cloud is not. In both species,
training on rows removes most of the neighbor penalty (91\,\% in maize and
85\,\% in sorghum, Sec.~\ref{sec:regime3}). The real scans
(Sec.~\ref{sec:res-field}) are not yet evaluated; until they are, the
one-camera transfer gap is our prediction for the sim-to-real penalty. A generator also contains unintended
regularities, such as the fixed leaf azimuth that gave a pose-aware baseline a
canonical orientation, and benchmark users should expect more of them.

\paragraph{What distillation does and does not buy.}
Distillation from a full-modality teacher improves every single-sensor student
and concentrates the gain on the weakest sources and the hardest plants
(Tab.~\ref{tab:anysubset}, Fig.~\ref{fig:views}B), and supervising all four
targets beats a parameter specialist at parameter regression. Distillation does
not make the model view-robust: the elevation curve moves down and keeps its
shape, so view robustness requires a different method, most plausibly treating
views as an explicit augmentation axis during pretraining. Distillation also
showed that an unpredictable field is not neutral: the three waviness parameters
of the first sorghum release stayed at zero skill and became more view-dependent
under distillation, and the final generator does not include them. A
parameterization meant to be regressed should be restricted to the fields that
the sensors can determine.

\paragraph{Scope and compute.}
The ablations use one seed per arm. The run-to-run floor comes from a single
near-duplicate pair (2\,\% in Chamfer, at most $0.011$ in $R^2$), and the maize
scaling differences, although significant under the paired bootstrap, could
still be run-to-run variance (Sec.~\ref{sec:res-repr}). Two generators and two
species-specific models are studied, and the 1{,}000-epoch headline model is
trained once per species. The procedural prior that makes inversion tractable
also limits it: a plant that the generator cannot express cannot be recovered,
and an organ class that the generator renders without variation (tassel droop in
maize is below zero for every encoder of this design) limits the data as much as
the model. Real scans and the \FieldProc{} release at
10K--100K plants remain \todo{in progress}.

\section{Conclusion}
\label{sec:conclusion}

\CropCode{} recovers a complete, editable procedural plant from the data of
whichever sensor a field campaign used, in one forward pass, by using the
generator for both the training data and the output representation. Two
findings from its ablations apply beyond plants: the generator's parameters are
valuable targets but harmful encoder inputs, and on slender geometry a
downstream probe, not Chamfer distance, should decide what a modality is worth.
A single camera's partial view requires refitting the readout and leaves a
measured transfer gap; the row-occluded and real-scan evaluations will show
whether the generator's prior covers the remaining difference to field data.
\todo{Sharpen once E10 lands.}
